\subsection{Математическое ожидание как интеграл Лебега-Стилтьеса.}
Пусть $\{x_i\}_{i = 1}^n$ -- возрастающая последовательность точек на числовой прямой и задана случайная величина $\xi$ на вероятностном пространстве $(\Omega, \mathcal{F}, \P)$. Рассмотрим случайную величину $S_n$, которая определяется как 
\[
    S_n(w) = \begin{cases}
        \sum\limits_{j = 1}^{n - 1} x_j\Ind\{w \mid x_j \leq \xi(w) < x_{j + 1} \}, & \xi(w) \in [x_1, x_n), \\
        0, \xi(w) < x_1 \vee \xi(w) \geq x_{n}.
    \end{cases} 
\]
Это простая функция, и существует такая последовательность $S_n$, что $S_n \uparrow \xi, n \ra \infty$. Поэтому, (для интегрируемых функций $\xi$) по теореме о Лебега мажорируемой сходимости сказать что
\[
    \Ex \xi = \int \xi d\P = \int \lim\limits_{n \ra \infty} S_n d\P = \lim\limits_{n \ra \infty} \Ex S_n.
\]
В тоже время
\[
    \Ex S_n = \sum\limits_{j = 1}^{n - 1} x_j(F_{\xi}(x_{j + 1}) - F_{\xi}(x_j)
\]
При устремлении $n \ra \infty$ 
\[
    \Ex \xi = \lim\limits_{n \ra \infty} \sum\limits_{j = 1}^{n - 1} x_j(F_\xi(x_{j + 1} - F_\xi(x_j)) = \int_{\R} xdF_\xi.
\]
Будем называть правую часть равенства <<интегралом Лебега-Стилтьеса>>. 
\begin{note}
    Интеграл Лебега-Стилтьеса -- это просто интеграл Лебега по мере, индуцированной на $(\R, \mathcal{B}(\R))$ при помощи случайной величины $\xi$, то есть по т.н. <<push-forward>> мере $\P_\xi(B) = \P(\xi^{-1}(B))$.
\end{note}
\subsubsection{Математическое ожидание дискретных и непрерывных случайных величин, примеры.}
Таким образом, для дискретных случайных величин (то есть случайных величин, мера которых сосредоточена лишь в счётном числе точек) мы можем считать математическое ожидание просто как 
\[
    \Ex \xi = \sum\limits_{j = 1}^{\infty} x_j \P(\xi(w) = x_j),
\]
где $\xi = \sum\limits_{j = 1}^{\infty} x_j\Ind_{A_j}$ и $\{A_j\}$ -- разбиение $\Omega$. \\
А математическое ожидание абсолютно непрерывных случайных величин (то есть величин таких, что их функция распределения имеет плотность относительно меры Лебега) мы можем считать математическое ожидание просто как 
\[
    \Ex \xi = \int_{\R} xdF_{\xi} = \int_{\R}x (F_{\xi})'dx = \int_{\R}xf_{\xi}dx.
\]
\begin{note}
    Олег Геннадьевич на лекции показывает <<корректность определения>> интеграла Лебега-Стилтьеса для математического ожидания сложной функции, то есть если $\eta = \phi(\xi)$, то
    \[
        \int_{\R}\phi(x)dF_{\xi} = \int_{\R} ydF_{\eta}(y).
    \]
    Кажется, в силу теоремы об абстрактной замене переменной в интеграле Лебега это действие достаточно бессмысленное: мы просто делаем второй push-forward $\xi$, то есть <<толкаем>> $\P_\xi$ при помощи $\phi$ и определяем $\P_\eta$. При таком подходе (то есть понимании того, что интеграл Лебега-Стилтьеса -- интеграл по мере, порождённой функцией распределения -- есть интеграл по push-forward мере случайной величины) эта проверка не нужна.
\end{note}
\subsubsection{Маргинальные распределения.}
\begin{definition}
    Пусть $\vec{\xi}$ -- случайный вектор, а $\vec{\xi}'$ -- его подвектор. Тогда распределение $\vec{\xi}'$ называется маргинальным.
\end{definition}
Функция распределения $\vec{\xi}'$ получается переходом к пределу (в $+\infty$) в функции распределения $\vec{\xi}$ по ненужным координатам. Если же $\vec{\xi}$ имеет плотность $f_{\vec{\xi}}$, то плотность маргинального распределения получается интегрированием по ненужным координатам.
\begin{note}[прим. техающего]
    Или я слепой, или этого не было в лекциях (посвящённых математическому ожиданию). Да и этот подпункт сюда тоже как-то не очень вписывается.
\end{note}